% Input after extension.tex.  Requires amsmath, amssymb, amsthm.
\subsection{从超临界侧趋近临界值的精确渐近}\label{subsec:critical-pole}

仍令 $m\ge2$、$a_i>0$，并采用 \textup{(W.1)} 的水位函数。记
\[
 \lambda_c=P_{\boldsymbol a}(1)>0,
 \qquad \sum_{i=1}^m\left(1-\frac{a_i}{\rho_c}\right)_+=1.
 \tag{CP.1}
\]
后式唯一确定 $\rho_c>0$。对 $\lambda>\lambda_c$，令
\[
 \beta(\lambda)=1-s_{\boldsymbol a}(\lambda).
 \tag{CP.2}
\]
记中心极大算子的弱型范数为
\[
 W_\mu=\sup_{\substack{0\ne f\in L^1(\mu)\\ t>0}}
 \frac{t\,\mu\{M_{c,\mu}f>t\}}{\|f\|_{L^1(\mu)}}.
 \tag{CP.3}
\]
中心方体的统一弱型界给出 $W_\mu\le2^n$；取有界正质量开方体的指标并令水平趋于一，得 $W_\mu\ge1$。

\begin{corollary}[全权族的全阈值公式]
对 $0<\lambda<1$，在非零非负局部可积权族中，
\[
\sup_w C_{c,w}^{\boldsymbol a}(\lambda)=
\begin{cases}
\infty,&0<\lambda\le\lambda_c,\\
1+2^n\dfrac{s_{\boldsymbol a}(\lambda)}{1-s_{\boldsymbol a}(\lambda)},&\lambda_c<\lambda<1.
\end{cases}
\tag{PH.1}
\]
\end{corollary}
\begin{proof}
临界及以下用定理 \ref{thm:general-critical-round11}。超临界时由定理 \ref{thm:weighted-transfer} 化为线性常数，再使用定理 \ref{endpoint} 的全权族精确公式。
\end{proof}

\begin{lemma}[小阈值指标函数的弱型范数]\label{lem:small-level-weak-norm}
若 $\mu$ 为非零、无原子、局部有限的 Radon 测度，则
\[
 \lim_{\beta\downarrow0}\beta C^{\rm int}_{c,\mu}(\beta)=W_\mu.
 \tag{CP.4}
\]
\end{lemma}

\begin{proof}
在定理 \ref{thm:two-measure-weak} 中取 $K=[-1,1]^n$、$\nu=\mu$，即得结论。
\end{proof}


\begin{theorem}[临界上侧的精确渐近]
令 $I_c=\{i:a_i<\rho_c\}$，$q=\#I_c$，
$A=\sum_{i\in I_c}a_i$。则 $q\ge2$、$\rho_c=A/(q-1)$，并且存在
$\beta_0>0$，使得对 $0\le\beta<\beta_0$，
\[
 P_{\boldsymbol a}(1-\beta)
 =\lambda_c\left(1+\frac{\beta}{q-1}\right)^A.
 \tag{CP.5}
\]
特别地，即使 $\rho_c$ 等于某个指数（活动集在 $s=1$ 切换），仍有左导数
\[
 P'_{\boldsymbol a}(1-)=-\rho_c\lambda_c,
 \qquad
 \beta(\lambda)\sim\frac{\lambda-\lambda_c}{\rho_c\lambda_c}
 \quad(\lambda\downarrow\lambda_c).
 \tag{CP.6}
\]

设 $\mu$ 为非零、无原子、局部有限的 Radon 测度。对 $\lambda_c<\lambda<1$，由定理
\ref{thm:weighted-transfer} 的精确传递，
\[
 \lim_{\lambda\downarrow\lambda_c}(\lambda-\lambda_c)
 C^{\boldsymbol a}_{c,\mu}(\lambda)=\rho_c\lambda_cW_\mu.
 \tag{CP.7}
\]
特别地，$C^{\boldsymbol a}_{c,\mu}(\lambda)
\asymp_{\mu,\boldsymbol a}(\lambda-\lambda_c)^{-1}$。

若 $d\mu=w\,dx$，并对所有非零非负 $w\in L^1_{\rm loc}$ 取上确界，
则方体族的已有精确公式给出
\[
 \lim_{\lambda\downarrow\lambda_c}
 (\lambda-\lambda_c)
 \sup_w C^{\boldsymbol a}_{c,w}(\lambda)
 =2^n\rho_c\lambda_c.
 \tag{CP.8}
\]
在原等指数特例 $a_1=\cdots=a_m=1$，右端为
$2^n(1-1/m)^{m-1}$。
\end{theorem}

\begin{proof}
由 \textup{(CP.1)}，在 $I_c$ 上求和得
\[
 1=q-\frac A{\rho_c},\qquad \rho_c=\frac A{q-1}.
 \tag{CP.9}
\]
由于每个正缺陷严格小于 $1$ 而总和为 $1$，有 $q\ge2$。当 $\beta>0$
充分小时，$A/(q-1+\beta)<\rho_c$，但仍严格大于每个
$i\in I_c$ 的 $a_i$。故水位引理 \ref{lem:water-filling} 表明，在
$s=1-\beta$，恰有 $I_c$ 中坐标活跃，水位为
\[
 \rho(1-\beta)=\frac A{q-1+\beta}.
\]
代回目标函数即得 \textup{(CP.5)}。这也说明当某些 $a_i=\rho_c$ 时，
这些坐标从左侧保持零缺陷，因而没有遗漏切换项。对 \textup{(CP.5)} 求导
（注意 $s=1-\beta$）便得 \textup{(CP.6)} 的第一式；再由一阶泰勒展开
得到第二式。

式 \textup{(CP.7)} 由 \textup{(CP.6)}、引理 \ref{lem:small-level-weak-norm} 以及精确传递立得。

最后，已有全权族恒等式写为
\[
 \sup_w C^{\boldsymbol a}_{c,w}(\lambda)
 =1+2^n\frac{s_{\boldsymbol a}(\lambda)}{1-s_{\boldsymbol a}(\lambda)}
 =1+2^n\frac{1-\beta(\lambda)}{\beta(\lambda)}.
\]
将它与 \textup{(CP.6)} 相乘并取极限，得到 \textup{(CP.8)}。
\end{proof}

\begin{remark}[一维 Lebesgue 情形]
当 $n=1$、$d\mu=dx$ 时，Melas 计算出中心 Hardy--Littlewood 极大算子的
弱型范数为 $W_{dx}=(11+\sqrt{61})/12$。因此此时 \textup{(CP.7)} 给出
固定 Lebesgue 测度的显式临界首项系数 $\rho_c\lambda_c(11+\sqrt{61})/12$
\cite{Melas2003}。例如 $m=2$、等指数时，固定 Lebesgue 测度的临界首项系数为 $(11+\sqrt{61})/24$，而对全部权取上确界的临界首项系数为 $1$。这两个量具有不同的量词。
\end{remark}
